\documentclass{article}
\usepackage{amsmath, amssymb}
\usepackage{amsthm}

\author{Michael Naumov}
\date{}
\setlength{\parindent}{0pt}
\setlength{\parskip}{0.5em}


\title{1.2.4\mbox{-}26}

\begin{document}
\maketitle

\( p \geq 3 \)

\subsubsection*{1}

If \( p \backslash a \implies a \equiv 0\ (\mathrm{modulo}\ p) \implies b = a \cdot a^{\frac{p-3}{2}} \equiv 0 \cdot a^{\frac{p-3}{2}} = 0 \ (\mathrm{modulo}\ p) \implies b \mod p = 0 \)

\subsubsection*{2}

If \( \lnot(p \backslash a) \)

Consider \( (b-1)(b+1)=a^{p-1}-1 \equiv 0\ (\mathrm{modulo}\ p) \)

\paragraph*{2.1}

If \( b - 1 \equiv 0\ (\mathrm{modulo}\ p) \implies b \mod p = 1 \)

\paragraph*{2.2}

If \( b - 1 \not\equiv 0\ (\mathrm{modulo}\ p) \implies (b - 1) \perp p \)

\( \implies b + 1 \equiv 0\ (\mathrm{modulo}\ p) \implies b \mod p = p - 1 \)


\end{document}
